% Simulazione d'esame: prova del 14/09/2021, esercizio 1 (B65-E01) -- dati iniziali isentropici.
% Stessi numeri e figure dello svolgimento di riferimento:
% svolgimenti/2021-09-14_B65-E01_dati_iniziali_isentropici (calcoli.py e figure/).
% Compilare con: bash simulazione_esame/compila.sh 2021-09-14_B65-E01_dati_iniziali_isentropici

% Preambolo comune delle simulazioni d'esame (vedi GUIDA.md).
% Ogni esercizio: % Preambolo comune delle simulazioni d'esame (vedi GUIDA.md).
% Ogni esercizio: % Preambolo comune delle simulazioni d'esame (vedi GUIDA.md).
% Ogni esercizio: \input{../preambolo} e \graphicspath verso le figure dello
% svolgimento di riferimento; si compila con compila.sh (svolgimento.pdf, A4).

\documentclass[11pt]{article}

\usepackage[T1]{fontenc}
\usepackage[utf8]{inputenc}
\usepackage[italian]{babel}
\usepackage{amsmath,amssymb}
\usepackage{array}
\usepackage{siunitx}
\usepackage{booktabs}
\usepackage{graphicx}
\usepackage{float}
\usepackage{caption}
\usepackage{enumitem}
\usepackage{hyperref}
\usepackage[a4paper,margin=2.5cm]{geometry}

\hypersetup{hidelinks}
\setlength{\parindent}{0pt}
\setlength{\parskip}{4pt}
\allowdisplaybreaks
\sisetup{output-decimal-marker={,}, group-separator={\,}, per-mode=symbol,
         exponent-product=\cdot, range-phrase={\text{--}}}

% Notazione del formulario (come template/svolgimento.tex).
\newcommand{\Cp}{C^{+}}
\newcommand{\Cm}{C^{-}}
\newcommand{\Mr}{M_{r}}
\newcommand{\Ms}{M_{s}}

% Separatore fra relazioni nella stessa formula: \begin{gather*} A \ee B \end{gather*}.
\newcommand{\ee}{,\qquad}

% Titolo dell'esercizio: \esercizio{Prova del GG/MM/AAAA -- Esercizio n}{argomento}.
\newcommand{\esercizio}[2]{%
  \title{#1\\[2pt]\large #2}\author{}\date{}%
  \hypersetup{pdftitle={#1: #2}}}

% Figura centrata: \figura[frazione della larghezza]{file}{didascalia\label{...}}.
\newcommand{\figura}[3][0.7]{%
  \begin{figure}[H]\centering\includegraphics[width=#1\textwidth]{#2}\caption{#3}\end{figure}}

% Testo a sinistra e figura (anche un dettaglio ingrandito) accanto, per tenere la figura vicino al punto che
% si discute: \accanto[larghezza della figura]{file}{didascalia\label{...}}{testo}.
\newcommand{\accanto}[4][0.40]{%
  \par\noindent
  \begin{minipage}[t]{\dimexpr\textwidth-#1\textwidth-0.03\textwidth\relax}\vspace{0pt}#4\end{minipage}\hfill
  \begin{minipage}[t]{#1\textwidth}\vspace{0pt}\centering\includegraphics[width=\linewidth]{#2}%
  \captionof{figure}{#3}\end{minipage}\par\medskip}
 e \graphicspath verso le figure dello
% svolgimento di riferimento; si compila con compila.sh (svolgimento.pdf, A4).

\documentclass[11pt]{article}

\usepackage[T1]{fontenc}
\usepackage[utf8]{inputenc}
\usepackage[italian]{babel}
\usepackage{amsmath,amssymb}
\usepackage{array}
\usepackage{siunitx}
\usepackage{booktabs}
\usepackage{graphicx}
\usepackage{float}
\usepackage{caption}
\usepackage{enumitem}
\usepackage{hyperref}
\usepackage[a4paper,margin=2.5cm]{geometry}

\hypersetup{hidelinks}
\setlength{\parindent}{0pt}
\setlength{\parskip}{4pt}
\allowdisplaybreaks
\sisetup{output-decimal-marker={,}, group-separator={\,}, per-mode=symbol,
         exponent-product=\cdot, range-phrase={\text{--}}}

% Notazione del formulario (come template/svolgimento.tex).
\newcommand{\Cp}{C^{+}}
\newcommand{\Cm}{C^{-}}
\newcommand{\Mr}{M_{r}}
\newcommand{\Ms}{M_{s}}

% Separatore fra relazioni nella stessa formula: \begin{gather*} A \ee B \end{gather*}.
\newcommand{\ee}{,\qquad}

% Titolo dell'esercizio: \esercizio{Prova del GG/MM/AAAA -- Esercizio n}{argomento}.
\newcommand{\esercizio}[2]{%
  \title{#1\\[2pt]\large #2}\author{}\date{}%
  \hypersetup{pdftitle={#1: #2}}}

% Figura centrata: \figura[frazione della larghezza]{file}{didascalia\label{...}}.
\newcommand{\figura}[3][0.7]{%
  \begin{figure}[H]\centering\includegraphics[width=#1\textwidth]{#2}\caption{#3}\end{figure}}

% Testo a sinistra e figura (anche un dettaglio ingrandito) accanto, per tenere la figura vicino al punto che
% si discute: \accanto[larghezza della figura]{file}{didascalia\label{...}}{testo}.
\newcommand{\accanto}[4][0.40]{%
  \par\noindent
  \begin{minipage}[t]{\dimexpr\textwidth-#1\textwidth-0.03\textwidth\relax}\vspace{0pt}#4\end{minipage}\hfill
  \begin{minipage}[t]{#1\textwidth}\vspace{0pt}\centering\includegraphics[width=\linewidth]{#2}%
  \captionof{figure}{#3}\end{minipage}\par\medskip}
 e \graphicspath verso le figure dello
% svolgimento di riferimento; si compila con compila.sh (svolgimento.pdf, A4).

\documentclass[11pt]{article}

\usepackage[T1]{fontenc}
\usepackage[utf8]{inputenc}
\usepackage[italian]{babel}
\usepackage{amsmath,amssymb}
\usepackage{array}
\usepackage{siunitx}
\usepackage{booktabs}
\usepackage{graphicx}
\usepackage{float}
\usepackage{caption}
\usepackage{enumitem}
\usepackage{hyperref}
\usepackage[a4paper,margin=2.5cm]{geometry}

\hypersetup{hidelinks}
\setlength{\parindent}{0pt}
\setlength{\parskip}{4pt}
\allowdisplaybreaks
\sisetup{output-decimal-marker={,}, group-separator={\,}, per-mode=symbol,
         exponent-product=\cdot, range-phrase={\text{--}}}

% Notazione del formulario (come template/svolgimento.tex).
\newcommand{\Cp}{C^{+}}
\newcommand{\Cm}{C^{-}}
\newcommand{\Mr}{M_{r}}
\newcommand{\Ms}{M_{s}}

% Separatore fra relazioni nella stessa formula: \begin{gather*} A \ee B \end{gather*}.
\newcommand{\ee}{,\qquad}

% Titolo dell'esercizio: \esercizio{Prova del GG/MM/AAAA -- Esercizio n}{argomento}.
\newcommand{\esercizio}[2]{%
  \title{#1\\[2pt]\large #2}\author{}\date{}%
  \hypersetup{pdftitle={#1: #2}}}

% Figura centrata: \figura[frazione della larghezza]{file}{didascalia\label{...}}.
\newcommand{\figura}[3][0.7]{%
  \begin{figure}[H]\centering\includegraphics[width=#1\textwidth]{#2}\caption{#3}\end{figure}}

% Testo a sinistra e figura (anche un dettaglio ingrandito) accanto, per tenere la figura vicino al punto che
% si discute: \accanto[larghezza della figura]{file}{didascalia\label{...}}{testo}.
\newcommand{\accanto}[4][0.40]{%
  \par\noindent
  \begin{minipage}[t]{\dimexpr\textwidth-#1\textwidth-0.03\textwidth\relax}\vspace{0pt}#4\end{minipage}\hfill
  \begin{minipage}[t]{#1\textwidth}\vspace{0pt}\centering\includegraphics[width=\linewidth]{#2}%
  \captionof{figure}{#3}\end{minipage}\par\medskip}

\graphicspath{{../../svolgimenti/2021-09-14_B65-E01_dati_iniziali_isentropici/figure/}}

\newcommand{\uma}{u-a}
\newcommand{\upa}{u+a}

\esercizio{Prova del 14/09/2021 -- Esercizio 1}{Dati iniziali isentropici: espansione e compressione centrata}

\begin{document}
\maketitle

\emph{Traccia.} Calcolare la soluzione per $t\ge0$ del problema 1D non stazionario con condizioni iniziali
isentropiche: per $x\le0$ $p=1$, $u=0$, $a=\num{1.522}$; per $0\le x\le1$ $u=\num{2.61}-\num{1.856}\,x$,
$a=1+\num{0.4995}\,x$; per $x\ge1$ $u=\num{0.754}$, $a=\num{1.4995}$. Graficare $p(x)$ e $M(x)$ a
$t=\num{0.4246}+\varepsilon$, stimare la soluzione asintotica, riportare la soluzione sui piani $x$--$t$ e $p$--$u$.

\section*{I dati e il modello}
I dati sono adimensionali e il gas non è indicato. Assumo $\gamma=\num{1.4}$, $\delta=\num{0.2}$, e non a caso:
è l'unico valore per cui $a+\delta u$ vale lo stesso ai due lati del salto in $x=0$ ($a_2+\delta u_2$ deve valere
$a_1$, con $u_1=0$), cioè per cui quel salto è un'onda semplice; con un altro $\gamma$ in $x=0$ nascerebbero tre
onde e il dato \num{1.522} sarebbe casuale. ``Isentropiche'' lo leggo come entropia uniforme: la pressione è data
solo a sinistra e altrove segue da $a$ con l'isentropica.

Il profilo è continuo in $x=1$ (la legge lineare dà proprio $u=\num{0.754}$, $a=\num{1.4995}$) e salta in $x=0$.
Chiamo 1 lo stato $x\le0$ ($u=0$, $a=\num{1.522}$), 2 lo stato in $x=0^+$ ($u=\num{2.61}$, $a=1$), 4 lo stato
$x\ge1$ ($u=\num{0.754}$, $a=\num{1.4995}$).

Uso le equazioni di Eulero 1D instazionarie per un gas perfetto non viscoso, in un dominio illimitato senza
contorni. Finché non nasce un urto il flusso è omoentropico: valgono gli invarianti di Riemann, $R_1=a-\delta u$
lungo le $\uma$ e $R_2=a+\delta u$ lungo le $\upa$, e $p/p_1=(a/a_1)^{2\gamma/(\gamma-1)}$ ovunque. Così
\[
p_2=\num{0.0529},\quad M_2=\num{2.61};\qquad p_4=\num{0.901},\quad M_4=\num{0.503}.
\]
Dopo la nascita dell'urto l'entropia non è più uniforme: gli invarianti valgono solo a tratti e lo stato dietro
l'urto lo trovo con Rankine--Hugoniot.

\section*{Quali onde ci sono: Mach e invarianti}
Il Mach con segno: $M_1=0$, in $x=0^+$ la corrente è \emph{supersonica} verso destra ($M_2=\num{2.61}$), nel tratto
lineare scende fino a $M_4=\num{0.503}$. L'entropia è uniforme per ipotesi. Per riconoscere le famiglie attive
calcolo gli invarianti:
\begin{table}[H]
  \centering
  \begin{tabular}{lll}
    \toprule
    & $R_1=a-\delta u$ & $R_2=a+\delta u$ \\
    \midrule
    stato 1 & \num{1.522} & \num{1.522} \\
    $0<x<1$ & $\num{0.478}+\num{0.871}\,x$ & $\num{1.522}+\num{0.128}\,x$ \\
    stato 4 & \num{1.349} & \num{1.650} \\
    \bottomrule
  \end{tabular}
\end{table}

\textbf{In $x=0$} $R_2$ è continuo e salta solo $R_1$: c'è \emph{solo un'onda $\uma$}, perché un'onda $\uma$ non
cambia $R_2$, trasportato dalle $\upa$ che la attraversano. La velocità caratteristica passa da $u_1-a_1=-\num{1.522}$
a $u_2-a_2=\num{1.61}$: cresce verso destra, le caratteristiche divergono, quindi è un'\textbf{espansione centrata
$\uma$} nell'origine (un salto di espansione non sarebbe ammissibile e si apre in ventaglio). Torna con la regola del
verso: l'espansione si propaga, rispetto al gas, verso sinistra e richiama il gas verso destra, da $u=0$ a
$u=\num{2.61}$. Testa e coda hanno segno opposto: l'espansione è \textbf{transonica}, con il punto sonico sulla
caratteristica verticale $x=0$. Nel ventaglio $\xi=x/t=u-a$ e $a+\delta u=R_{2,1}$, quindi
\[
u=\frac{\xi+R_{2,1}}{1+\delta}\ee a=\frac{R_{2,1}-\delta\,\xi}{1+\delta}\ee
-\num{1.522}\le\xi\le\num{1.61};
\]
in $\xi=0$ trovo $u=a=R_{2,1}/(1+\delta)=\num{1.268}$, $p=\num{0.279}$. Controllo il vuoto: nel ventaglio
$a\ge1>0$, nessun problema.

\textbf{Nel tratto $0<x<1$} variano entrambi gli invarianti: ci sono \emph{tutte e due le famiglie}, e il tratto è
all'inizio una zona d'interazione. Per capire che cosa diventa ciascuna famiglia scrivo le velocità caratteristiche
in funzione degli invarianti, da $u=(R_2-R_1)/(2\delta)$ e $a=(R_1+R_2)/2$:
\[
u-a=\frac{(1-\delta)R_2-(1+\delta)R_1}{2\delta}=2R_2-3R_1\ee
u+a=\frac{(1+\delta)R_2-(1-\delta)R_1}{2\delta}=3R_2-2R_1 .
\]
\begin{itemize}
\item Le $\uma$ portano $R_1$, che cresce con $x$. Le $\upa$ che arrivano da sinistra (dallo stato 1 e dal
  ventaglio) portano tutte $R_2=\num{1.522}$, quindi, appena le famiglie si separano, le $\uma$ viaggiano a
  $2R_{2,1}-3R_1(x_0)$, velocità che decresce con il punto di partenza: è una \textbf{compressione $\uma$}, che
  formerà un urto. Già a $t=0$ $u-a=\num{1.61}-\num{2.36}\,x$ decresce.
\item Le $\upa$ portano $R_2$, che cresce con $x$. Dopo la separazione attraversano lo stato 4, vedono
  $R_1=R_{1,4}$ e viaggiano a $3R_2(x_0)-2R_{1,4}$, velocità \emph{crescente}: è un'\textbf{espansione $\upa$} non
  centrata verso destra. A $t=0$ $u+a=\num{3.61}-\num{1.36}\,x$ decresce solo perché vi è sovrapposta la
  compressione $\uma$: è un effetto dell'interazione, non una compressione $\upa$.
\end{itemize}

\section*{La compressione $\uma$ è centrata}
Seguo l'istante dato dalla traccia, cioè la formula dell'onda semplice con le $\uma$ rette uscenti dai punti
iniziali, $x=x_0+(u-a)^0(x_0)\,t$ (il tratto non è esattamente un'onda semplice: lo discuto in fondo). Due
caratteristiche vicine si incontrano quando $1+(u-a)^{0\prime}\,t=0$:
\[
t^*=\frac{1}{|(u-a)^{0\prime}|}=\num{0.4246}\ee x^*=(u_2-a_2)\,t^*=\num{0.684}.
\]
Con $(u-a)^0$ lineare, $x_0+(u-a)^0(x_0)\,t^*=x^*$ per \emph{ogni} $x_0\in[0,1]$: tutte le $\uma$ della
compressione arrivano nello stesso punto $C=(x^*,t^*)$, compresa la prima, che è la coda del ventaglio. Non c'è
coalescenza graduale: in $C$ nasce subito un urto di ampiezza finita. Per la famiglia $\upa$ la stessa formula
darebbe $1/\num{1.36}=\num{0.737}$, ma non ha significato perché quella famiglia, dopo la separazione, è
un'espansione.

\section*{Lo stato fra le due famiglie}
Fra la compressione $\uma$ e l'onda $\upa$, una volta separate, c'è uno stato uniforme 3, e per trovarlo non serve
risolvere la zona d'interazione: le $\upa$ che vi arrivano portano $R_2$ dello stato 1, le $\uma$ portano $R_1$
dello stato 4.
\[
a_3=\frac{R_{2,1}+R_{1,4}}{2}=\num{1.435}\ee u_3=\frac{R_{2,1}-R_{1,4}}{2\delta}=\num{0.433}\ee
p_3=\num{0.663}\ee M_3=\num{0.302}.
\]
Pendenze: $u_3-a_3=-\num{1.00}$, $u_3+a_3=\num{1.87}$; nello stato 4 $u_4-a_4=-\num{0.746}$, $u_4+a_4=\num{2.25}$.

Devo controllare che a $t^*$ le famiglie siano davvero separate, altrimenti a destra di $C$ non ci sarebbe lo
stato 3. La separazione avviene nel punto $S$ in cui la $\upa$ da $x=0$ (pendenza da $u_2+a_2=\num{3.61}$ a
$u_3+a_3$) incontra la $\uma$ da $x=1$ (da $u_4-a_4$ a $u_3-a_3$). Con le pendenze medie, come nel reticolo delle
caratteristiche, trovo $t_S\approx\num{0.28}$, $x_S\approx\num{0.76}$: $t_S<t^*$, verificato. La zona a due
famiglie è il triangolo di base $0\le x\le1$ chiuso da $S$.

A $t^*$ l'onda $\upa$ occupa allora $x_S+(u_3+a_3)(t^*-t_S)\approx\num{1.03}\le x\le1+(u_4+a_4)t^*=\num{1.96}$:
la testa è esatta, perché viaggia nello stato 4 uniforme, la coda è stimata. Dopo $S$ le sue caratteristiche
viaggiano a $3R_2(x_0)-2R_{1,4}$, da \num{1.87} a \num{2.25}.

\section*{Il centro di compressione: un problema di Riemann}
Il centro di una compressione centrata è un punto singolare, quindi in $C$ risolvo un problema di Riemann fra lo
stato 2 che arriva da sinistra (sulla coda del ventaglio) e lo stato 3 a destra. Le incognite sono $\rho,u,p$ degli
stati 5 (dietro l'urto) e 6 (fra contatto e onda destra) e la velocità $W$ dell'urto: 7. Le equazioni sono 3 di
Rankine--Hugoniot, 2 sul contatto ($p_5=p_6$, $u_5=u_6$) e 2 attraverso l'onda destra (invariante $R_1$ e stessa
entropia di 3): 7, il problema è chiuso.

Primo tentativo isentropico sulla velocità: 2 e 3 stanno sulla stessa isentropica $\uma$ (stesso $R_2$), quindi
$u^{(0)}=u_3=\num{0.433}$. Con $u^{(0)}<u_2$ a sinistra c'è un urto, con $u^{(0)}\le u_3$ a destra
un'espansione. Le pressioni al primo tentativo sono $p_2^{(0)}=\num{0.528}$ con l'urto e $p_3^{(0)}=\num{0.663}$
con l'isentropica: la differenza è grande perché l'urto è forte. Mi aspettavo questo: la curva di
Rankine--Hugoniot sta sotto l'isentropica, quindi l'urto $\uma$ da solo non arriva allo stato 3 e servono anche
un contatto e un'espansione $\upa$.

Per l'urto ricavo il Mach relativo dal salto di velocità: dalle Rankine--Hugoniot nel riferimento dell'urto
$u_2-u=\tfrac{2a_2}{\gamma+1}\left(\Mr-1/\Mr\right)$, cioè $\Mr^2-2X\Mr-1=0$, e
\[
\Mr=X+\sqrt{1+X^2}\ee X=\frac{\gamma+1}{4}\,\frac{u_2-u}{a_2}\ee
\frac{p_5}{p_2}=1+\frac{2\gamma}{\gamma+1}\left(\Mr^2-1\right).
\]
Per l'espansione $a_6-\delta u_6=a_3-\delta u_3$ e $p_6=p_3(a_6/a_3)^{2\gamma/(\gamma-1)}$. Aggiorno $u$ finché
$|p_5-p_6|$ è trascurabile: dopo tre passi lo scarto è $\sim10^{-6}$, fisicamente nullo, e
\[
u^*=\num{0.300}\ee p^*=\num{0.582}.
\]
\begin{itemize}
\item \textbf{Urto $\uma$} (2$\to$5): $\Ms=\num{3.10}$, $W=u_2-a_2\Ms=-\num{0.485}$, $p_5/p_2=\num{11.0}$;
  \[
  \frac{\rho_5}{\rho_2}=\frac{(\gamma+1)\Ms^2}{(\gamma-1)\Ms^2+2}=\num{3.94}\ee
  a_5=\sqrt{\gamma p_5/\rho_5}=\num{1.67}\ee M_5=\num{0.180}.
  \]
\item \textbf{Contatto} a velocità $u^*=\num{0.300}$.
\item \textbf{Espansione $\upa$} (6$\to$3) centrata in $C$: $a_6=a_3-\delta(u_3-u^*)=\num{1.41}$,
  $M_6=\num{0.213}$; coda $u^*+a_6=\num{1.71}$, testa $u_3+a_3=\num{1.87}$.
\end{itemize}
Controlli subito: $\Ms=\num{3.10}>1$ e $s_5/s_2=\num{1.61}>1$, l'urto è ammissibile; $u^*<u_2$ (urto a sinistra) e
$u^*<u_3$ (espansione a destra), la natura delle onde è quella ipotizzata; condizione di Lax
$u_2-a_2=\num{1.61}>W=-\num{0.485}>u_5-a_5=-\num{1.37}$, le $\uma$ entrano nell'urto da entrambi i lati.

\section*{I profili a $t^*+\varepsilon$}
A $t^*+\varepsilon$ urto, contatto ed espansione nati in $C$ occupano un tratto di ampiezza $O(\varepsilon)$: in
$x^*$ si vede un solo salto da 2 a 3 che contiene la struttura 2--5--6--3. Da sinistra:
\begin{itemize}
\item $x<-a_1t^*=-\num{0.646}$: stato 1, $p=1$, $M=0$;
\item $-\num{0.646}<x<\num{0.684}$: ventaglio, con $\xi=x/t^*$ nelle formule sopra; $M=1$ in $x=0$;
\item $x=x^*=\num{0.684}$: salto da $p_2=\num{0.0529}$, $M_2=\num{2.61}$ a $p_3=\num{0.663}$, $M_3=\num{0.302}$;
\item $\num{0.684}<x\lesssim\num{1.03}$: stato 3;
\item $\num{1.03}\lesssim x<\num{1.96}$: espansione $\upa$ da 3 a 4 ($p$ da \num{0.663} a \num{0.901}, $M$ da
  \num{0.302} a \num{0.503}), con andamento interno stimato;
\item $x>\num{1.96}$: stato 4.
\end{itemize}

\figura[0.78]{profili_tstar.pdf}{$p(x)$ in scala lineare e logaritmica e $M(x)$ a $t=0$ (tratteggio) e a
$t=t^*+\varepsilon$.\label{fig:profili}}

In fig.~\ref{fig:profili} il ventaglio ha portato la pressione da 1 a \num{0.053} con andamento continuo e il Mach
da 0 a \num{2.61}, passando per $M=1$ in $x=0$; la compressione, che a $t=0$ occupava $0<x<1$, si è ridotta al salto
verticale in $x^*$, dove $p$ risale di un ordine di grandezza e $M$ crolla a \num{0.3}. A destra ci sono il plateau
3 e la rampa dell'onda $\upa$ fino allo stato 4. Poiché la pressione copre più di un ordine di grandezza riporto
entrambe le scale: la lineare mostra il peso relativo dei salti, la logaritmica rende leggibile la parte bassa del
ventaglio e l'entità del salto in $x^*$ (oltre un fattore 10).

\figura[0.95]{struttura_salto.pdf}{Struttura interna del salto in $x^*$ nella variabile $(x-x^*)/(t-t^*)$: urto,
contatto (tratteggi) ed espansione $\upa$ (fascia grigia).\label{fig:salto}}

La fig.~\ref{fig:salto} ingrandisce quel salto: l'urto porta $p$ da \num{0.053} a \num{0.582}; il contatto non
cambia $p$ ma cambia $M$ (da \num{0.180} a \num{0.213}), perché lo stato 5, riscaldato dall'urto, ha velocità del
suono maggiore; l'espansione $\upa$ completa il passaggio a 3.

\section*{Il piano $x$--$t$}
\begin{table}[H]
  \centering
  \small
  \begin{tabular}{clp{0.5\linewidth}}
    \toprule
    Stato & Dove & Onde che lo delimitano \\
    \midrule
    1 & $x<-a_1t$ & testa del ventaglio $\uma$ \\
    2 & sulla retta $x=\num{1.61}\,t$ & coda del ventaglio e prima caratteristica della compressione (spessore nullo) \\
    3 & fra le due famiglie separate & compressione $\uma$ (poi urto) a sinistra, onda $\upa$ a destra \\
    4 & $x>1+(u_4+a_4)t$ & testa dell'onda $\upa$ \\
    5 & dopo $t^*$, fra urto e contatto & urto $\uma$, contatto \\
    6 & dopo $t^*$, fra contatto ed espansione & contatto, espansione $\upa$ centrata in $C$ \\
    \bottomrule
  \end{tabular}
\end{table}

\figura[0.82]{xt.pdf}{Piano $x$--$t$: $\uma$ in blu, $\upa$ in verde, urto in rosso, contatto tratteggiato; in grigio
la zona a due famiglie. Dopo $t^*$ urto e contatto sono qualitativi.\label{fig:xt}}

Nel piano $x$--$t$ (fig.~\ref{fig:xt}) le zone uniformi sono 1, 3 e 4; le onde semplici, con caratteristiche
rettilinee, sono il ventaglio centrato nell'origine, la compressione $\uma$ convergente in $C$ (la coda del
ventaglio coincide con la sua prima caratteristica) e l'onda $\upa$; la zona d'interazione a due famiglie è il
triangolo grigio chiuso da $S$. Le $\upa$ uscenti dal tratto lineare, piegate nella zona d'interazione, ne escono
divergenti (espansione). Da $C$ partono l'urto, che piega verso sinistra entrando nel ventaglio e tende alla
pendenza $u_\infty-a_\infty$, il contatto, che tende alla velocità $u_\infty$, e la stretta espansione $\upa$ fra 6
e 3. Le $\uma$ che vengono dallo stato 4 (tratteggiate) attraversano l'onda $\upa$ e lo stato 3 e raggiungeranno
anch'esse l'urto, da valle.

\section*{La soluzione asintotica}
Per $t\to\infty$ la regione iniziale di lunghezza 1 diventa un punto e conta solo il salto fra gli stati esterni:
risolvo il pseudo-Riemann fra 1 e 4. Hanno la stessa entropia, quindi le onde sono isentropiche e lo stato
centrale segue dagli invarianti, $a_\infty+\delta u_\infty=R_{2,1}$ e $a_\infty-\delta u_\infty=R_{1,4}$: sono le
stesse due condizioni dello stato 3, quindi lo stato asintotico è già noto,
\[
p_\infty=\num{0.663}\ee u_\infty=\num{0.433}\ee M_\infty=\num{0.302}.
\]
Con $u_1<u_\infty<u_4$ entrambe le onde sono espansioni (RCR); $a_\infty>0$ e non c'è vuoto, perché
$u_4-u_1\ll(a_1+a_4)/\delta$.
\begin{itemize}
\item espansione $\uma$ da 1: testa $-\num{1.52}$, coda $u_\infty-a_\infty=-\num{1.00}$, \emph{subsonica};
\item espansione $\upa$ verso 4: coda $u_\infty+a_\infty=\num{1.87}$, testa $u_4+a_4=\num{2.25}$;
\item al posto del contatto, uno strato di gas a entropia più alta (le particelle passate attraverso l'urto),
  trasportato con $u_\infty$: $p$ e $u$ sono continui, $\rho$ e $T$ no.
\end{itemize}
L'urto, raggiunto dalle $\uma$ del ventaglio, si indebolisce progressivamente e tende a un'onda di Mach lungo la
coda $dx/dt=u_\infty-a_\infty$.

\section*{Il piano $p$--$u$}
\figura[0.62]{pu_lin.pdf}{Piano $p$--$u$, scala lineare.\label{fig:pulin}}
\figura[0.62]{pu_log.pdf}{Piano $p$--$u$, scala logaritmica.\label{fig:pulog}}

Gli stati 1, 2 e 3 stanno sulla stessa isentropica $\uma$ ($R_2=\num{1.522}$): il ventaglio porta da 1 a 2, la
compressione riporterebbe da 2 a 3 lungo la stessa curva. Gli stati 3 e 4 stanno sull'isentropica $\upa$ di 4. I
dati iniziali del tratto lineare (punteggiato) non stanno su nessuna delle due curve: è il segno che sono presenti
entrambe le famiglie. In $C$ la curva di Rankine--Hugoniot da 2 sta sotto l'isentropica e taglia l'isentropica
$\upa$ in $5=6$, a velocità e pressione minori di 3: per questo dietro l'urto serve un'espansione. Il contatto non
si vede (5 e 6 coincidono, ma con impedenze $\rho a$ diverse, \num{0.49} e \num{0.58}). Lo stato asintotico è
l'intersezione delle isentropiche di 1 e di 4, cioè 3. La scala lineare (fig.~\ref{fig:pulin}) mostra bene la zona
di $C$; quella logaritmica (fig.~\ref{fig:pulog}) rende leggibile la regione di bassa pressione vicino a 2, dove la
curva di Rankine--Hugoniot e l'isentropica partono tangenti (stessa tangente e curvatura) e si separano solo per
salti forti.

\section*{Risultati e discussione}
\begin{table}[H]
  \centering
  \small
  \begin{tabular}{lp{6.4cm}p{3.2cm}}
    \toprule
    Grandezza & Valore & Esatto / stimato \\
    \midrule
    Ventaglio $\uma$ in $(0,0)$ & $-\num{1.522}\le x/t\le\num{1.61}$, transonico & esatto \\
    Formazione dell'urto & $t^*=\num{0.4246}$, $x^*=\num{0.684}$ & caratteristiche rette \\
    Stato 3 & $p=\num{0.663}$, $u=\num{0.433}$, $M=\num{0.302}$ & esatto \\
    Urto in $C$ & $\Ms=\num{3.10}$, $W=-\num{0.485}$, $p_5/p_2=\num{11.0}$ & esatto (dato $C$) \\
    Stati 5 / 6 & $p^*=\num{0.582}$, $u^*=\num{0.300}$, $M_5=\num{0.180}$, $M_6=\num{0.213}$ & esatto \\
    Onda $\upa$ a $t^*$ & $\num{1.03}\lesssim x<\num{1.96}$ & coda stimata, testa esatta \\
    Asintotico & $p_\infty=\num{0.663}$, $u_\infty=\num{0.433}$, $M_\infty=\num{0.302}$, RCR & stima \\
    \bottomrule
  \end{tabular}
\end{table}

\textbf{Istante di formazione.} Il tratto lineare non è un'onda semplice ($R_2$ varia di circa l'8\%): nella breve
zona d'interazione le $\uma$ si piegano un po' verso sinistra (quella da $x=1$ passa da $u_4-a_4=-\num{0.75}$ a
$u_3-a_3=-\num{1.00}$). Con le pendenze dopo la separazione, $u-a=2R_{2,1}-3R_1(x_0)$, il gradiente sarebbe
\num{2.61} e $t\approx\num{0.38}$: l'urto nasce un po' prima di $t^*$ e la coalescenza non è esattamente in un
punto. La struttura (urto forte nato quasi tutto insieme, contatto, espansione $\upa$) non cambia; ho usato
l'istante della traccia.

\textbf{Perché nasce un urto forte.} La corrente supersonica a $M=\num{2.61}$ della coda del ventaglio va a
sbattere contro il gas lento dello stato 3, come contro un pistone: la compressione centrata non ha il tempo di
formare l'urto per gradi, che nasce già con $\Ms\approx3$. A questa intensità l'approssimazione isentropica
dell'urto (lecita per $\Ms\lesssim\num{1.2}$) non vale: darebbe direttamente lo stato 3, mentre le
Rankine--Hugoniot danno $p^*=\num{0.582}$ e un'espansione $\upa$ non trascurabile. L'entropia prodotta
($s_5/s_2=\num{1.61}$) resta nelle particelle urtate, che si comportano come una bolla calda.

\textbf{Evoluzione e regime asintotico.} L'espansione ``vive per sempre'', la compressione degenera: l'urto viaggia
a $W<0$ dentro un ventaglio le cui $\uma$ (più veloci) lo raggiungono da monte, quindi vede a monte gas sempre più
lento e caldo e si indebolisce. Alla fine la parte supersonica del ventaglio è stata assorbita: l'espansione
asintotica verso sinistra è subsonica (coda a $M=\num{0.30}$) e il salto netto fra 1 e 4 è una doppia espansione
con lo stato centrale 3. Il pseudo-Riemann dà gli stati asintotici, non i tempi dell'interazione urto--ventaglio,
che si potrebbe seguire solo con le caratteristiche e le relazioni d'urto punto per punto.
\end{document}
